
\documentclass[12pt,fullpage]{article}

\usepackage{amsthm}
\usepackage{amssymb}
\usepackage[fleqn]{amsmath}
\usepackage{graphicx}
\usepackage{enumitem}


\newtheorem{theorem}{Theorem}[section]
\newtheorem{lemma}[theorem]{Lemma}
\newtheorem{proposition}[theorem]{Proposition}
\newtheorem{corollary}[theorem]{Corollary}


\theoremstyle{definition}
\newtheorem{question}[theorem]{Question}
\newtheorem{definition}[theorem]{Definition}
\newtheorem{example}[theorem]{Example}
\newtheorem{xca}[theorem]{Exercise}

\theoremstyle{remark}
\newtheorem{remark}[theorem]{Remark}




\newcommand{\A}{\mathcal{A}}
\newcommand{\B}{\mathcal{B}}
\newcommand{\C}{\mathcal{C}}
\newcommand{\G}{\mathcal{G}}
\newcommand{\Q}{{\mathbb{Q}}}
\newcommand{\M}{\mathcal{M}}
\newcommand{\Z}{{\mathbb{Z}}}
\renewcommand{\P}{\mathcal{P}}
\newcommand{\D}{\mathcal{D}}
\newcommand{\AND}{\, \wedge \,}
\newcommand{\OR}{\, \vee \,}

\renewcommand{\phi}{\varphi}

\newcommand{\N}{\mathbb{N}}
\newcommand{\AUT}{\mbox{\it AUT}}
\newcommand{\iso}{{\stackrel{\sim}{\rightarrow}}}
\renewcommand{\restriction}{\upharpoonright}
\newcommand{\bigdoublewedge}{%
  \mathop{
    \mathchoice{\bigwedge\mkern-15mu\bigwedge}
               {\bigwedge\mkern-12.5mu\bigwedge}
               {\bigwedge\mkern-12.5mu\bigwedge}
               {\bigwedge\mkern-11mu\bigwedge}
    }
}

\newcommand{\bigdoublevee}{%
  \mathop{
    \mathchoice{\bigvee\mkern-15mu\bigvee}
               {\bigvee\mkern-12.5mu\bigvee}
               {\bigvee\mkern-12.5mu\bigvee}
               {\bigvee\mkern-11mu\bigvee}
    }
}
\begin{document}


\title{Descriptive Aspects of Structures with Countable Automorphism Groups}
\author{Dariusz Kalociński and  Theodore A. Slaman}
 
\maketitle

\section{Question}

We consider infinite countable first-order structures, their automorphism groups and the descriptive aspects of their representations. A representation $\A$ of such a structure is an instance of its isomorphism type in which the universe or domain of the structure is the set of natural numbers $\N$ and the constants, functions and relations of the structure are elements of $\N$, multi-variable functions from $\N$ to $\N$ and multi-place relations on $\N$, respectively.   Similarly, an automorphism of $\A$ or an isomorphism from $\A$ to $\B$ is a bijective function from $\N$ to $\N$ that preserves the evaluation of all constants, functions and relations.  We let $\AUT(\A)$ denote the automorphism group of $\A$.


\begin{definition}\label{1.1}
  \begin{itemize}
  \item     $\A$ is \emph{computably $\AUT$-countable} if every automorphism of $\A$ is computable relative to $\A$.  
  \item  $\A$ is \emph{computably $\AUT$-countable on a cone} if there is a $C\subset\N$ such that whenever $\B$ is isomorphic to $\A$, every automorphism of $\B$ is computable relative to $\B\oplus C$.
  \end{itemize}
\end{definition}

\begin{definition}
  A function $f:\A\to\A$ is \emph{$\Sigma^{in}_1$-definable} in $\A$ relative to parameters
  $\bar{a}=(a_1,\dots,a_n)$ from $\A$ if there is a family of existential formulas $(\phi_i)_{i\in\N}$ such that for all $x$ and $y$ in $\A$, $f(x)=y$ if and only if there is an $i\in\N$ such that $\A\models\phi_i(\bar{a},x,y)$.
\end{definition}

\begin{question}\label{sec:question}\label{1.2}
  Is there a countable computably-represented structure $\A$ such that $\A$ is computably $\AUT$-countable on a cone and such that for any finite  set of parameters $\bar{a}$ in $\A$ there is an automorphism $\pi$ of $\A$ such that $\pi$ is not $\Sigma^{in}_1$-definable in $\A$ relative to the parameters $\bar{a}\cup \pi(\bar{a}),$ where $\pi(\bar{a})$ is the image of $\bar{a}$ under $\pi$.
\end{question}



\section{Background}

We note some examples.

\begin{example}
  \begin{itemize}
  \item Any rigid structure, having no nontrivial automorphism, is computably $\AUT$-countable, since the identity map is computable no matter how complicated the structure may be.

  \item The first order structure represented by the integers with the successor function is computably $\AUT$-countable, since every automorphism is a finite iteration of the successor or of its inverse.  
  \end{itemize}
\end{example}


We recall a fundamental model theoretic result.

\begin{theorem}[see~\protect{\cite[Cor.\ 4.2.5]{Hod1993}}]\label{2.2}
  Let $\A$ be a countable structure and suppose that $\AUT(\A)$ is countable.  There are finitely many parameters $\bar{a}\in\A$ such that every automorphism of $\A$ is determined by its action on $\bar{a}$.
\end{theorem}

\begin{theorem}[following~\protect{\cite{AshKniManSla89}}]\label{2.3}
  Let $\A$ be a countable structure.  $\A$ is computably $\AUT$-countable on a cone if and only if for each $\pi\in\AUT(\A)$, there are finitely many parameters $\bar{a}\in\A$ such that $\pi$ is $\Sigma^{in}_1$-definable in $\A$ relative to $\bar{a}$.
\end{theorem}

\cite{AshKniManSla89}\ showed that for a relation $R$ on $\A$, $R$ is computably enumerable on a cone if and only if there are finitely many parameters $\bar{a}\in\A$ and a family of existential formulas $(\phi_i)_{i\in\N}$ such that all $\bar{x}\subset\A$, $\bar{x}\in R$ if and only if there is an $i$ such that $\A\models\phi_i(\bar{a},x)$.  We obtain Theorem~\ref{2.3} by applying this result to the graphs of the automorphisms of $\A$.

When its hypotheses apply,  Theorem~\ref{2.3} gives a specific syntactic form such that if $\pi$ is an automorphism of $\A$, then there are finitely many parameters within $\A$ such that the action of $\pi$ is definable from those parameters.  In the result as stated, the parameters used to define $\pi$ depend on which $\pi$ is being defined.  This does not address the full content of Theorem~\ref{2.2}, which asserts that there are finitely many parameters such that every automorphism is determined by its action on those parameters.  In Question~\ref{1.2}, we ask whether there is a counter-example to what would have been a natural effectivization of Theorem~\ref{2.2}, that there would be one finite set of parameters from which every automorphism could be effectively defined.



\section{Solution}
\label{sec:solution}

\begin{theorem}\label{3.1}
  There is a computably-represented structure $\A$ such that $\A$ is computably $\AUT$-countable on a cone and such that for any finite set of parameters $\bar{a}$ there is an automorphism $\pi$ of $\A$ such that $\pi$ is not $\Sigma^{in}_1$-definable in $\A$ relative to $\bar{a}\cup\pi(\bar{a})$.
\end{theorem}


Typically, counter-examples to computability-theoretic statements are produced by diagonal constructions.  This is not the case here.  Rather, the structure exhibited in Theorem~\ref{3.1} comes from properties of elementary arithmetic, specifically those of $(\Q,+)$, the additive group on the rational numbers.  

As is well-known, $\pi$ is an automorphism of $(\Q,+)$ if and only if there is a non-zero rational $s$ such that for all $x$, $\pi(x)=sx$.  Equivalently, the automorphism group of $(\Q,+)$ is naturally isomorphic to the multiplicative group on the non-zero rational numbers, conventionally denoted as  $(\Q^*,\cdot)$.  Here is a quick reminder of the proof.  To verify the non-trivial implication, suppose that $\pi$ is an automorphism of $(\Q,+)$ and let $a/b$ be the reduced-form fraction that represents $\pi(1)$.  Consider a non-zero rational $x$ represented by $c/d$.  Then
\[
  \underbrace {x+\dots+x}_{\mbox{$d$-times}}=
  \underbrace {1+\dots+1}_{\mbox{$c$-times}}
\]
and so 
\[
  \underbrace  {\pi(x)+\dots+\pi(x)}_{\mbox{$d$-times}}=
  \underbrace {a/b+\dots+a/b}_{\mbox{$c$-times}}.
\]
It follows that $\pi(x)=\frac{ca/b}{d}=a/b\cdot c/d=a/b \cdot x$.  For the case when $x=0$, $\pi(0)=0=a/b\cdot 0$, as required to show for all $x$, $\pi:x\mapsto a/b \cdot x.$


Before giving the proof of Theorem~\ref{3.1}, we note that $(\Q,+)$ is not an example of the desired sort.  If $\pi:x\mapsto a/b\cdot x$, then $\pi(x)=y$ exactly when \[
  (\Q,+)\models  \underbrace {y+\dots+y}_{\mbox{$b$-times}}=
  \underbrace {x+\dots+x}_{\mbox{$a$-times}},
\]
which shows that the graph of $\pi$ is definable in $(\Q,+)$ without parameters by a formula without quantifiers.   We produce an example $\A$ to verify Theorem~\ref{3.1}, by rearranging the arithmetic of $(\Q,+)$.

\begin{proof}[Proof of Theorem~\ref{3.1}]
  The language of $\A$ consists of two unary predicates $N$ and $Q$, a binary relation symbol $<$ and a ternary relation symbol $R$.
  The structure $\A$ consists of
  \begin{enumerate}
  \item A copy of the rational numbers,  which constitutes the elements of $\A$ which satisfy $Q$ in $\A.$
  \item A copy of the natural numbers, disjoint from $Q$,  which constitutes the elements of $\A$ which satisfy $N$ in $\A.$
  \item In $\A$, $n<m$ if and only if $n$ and $m$ are in $N$ and $n$ is less than $m$ in the natural order.
  \item Fix a computable counting $S_0,S_1,\dots$ of the finite subsets of the non-zero rational numbers so that for each such finite set $S$  there are infinitely many $n$ such that $S_n=S.$  In $\A$, $R(n,x,y)$ holds if and only if $n\in N$ and there is an $s\in S_n$ such that $y=sx$, equivalently $y/x\in S_n$, for $x$ nonzero.
  \end{enumerate}

 To simplify notation, we identify $N$, $Q$, $<$ and $R$ with their interpretations in $\A$.
  
  \paragraph{Representing elements of $\AUT(\A)$.} We claim that, just as for $(\Q,+),$ the automorphisms of $\A$ are represented by scalar multiplication by non-zero rational constants.  By inspection, if $s$ is a nonzero rational number, then the map on $Q$ given by $x\mapsto sx$ induces an automorphism of $\A$ (which is the identity on $N$).  For the converse, suppose that $\pi$ is an automorphism of $\A.$  We claim that there is a non-zero rational number $s$ such that $\pi$ is induced by the map $x\mapsto sx$ on $Q$.

To prove the claim, begin by noting that since every element of $N$ is definable, $\pi$ fixes the elements of $N$.  Consequently, $\pi$ fixes the sets $S_n$ which define the relation $R$.

Next,  we show that addition on $Q$ can be evaluated using the first order structure of $\A$.  Suppose that $a$ and $b$ are elements of $Q$.  If $a=0$, an $\A$-definable condition, then $a+b=b.$    Otherwise, begin to evaluate $a+b$ by letting $n$ be least such that $S_n$ is a singleton and $R(n,a,b)$.  In other words, $S_n=\{s\}$ and $b=sa$.  Next, let $m$ be least such that $S_m=\{1+s\}$.  Finally, evaluate $a+b$ by noting that $a+b=c$ if and only if $R(m,a,c)$.  Thus, $\pi$ fixes the additive structure on $Q$.  As noted earlier, the automorphisms of $(\Q,+)$ are exactly the functions $x\mapsto sx$, for $s$ a non-zero rational number, and the claim follows.  

\paragraph{Computable $\AUT$-countability on cone.}  Suppose that $\A^*$ is computable from $X$ and $f:\A^*\iso \A$.  Then, an isomorphism $\pi^*$ of $\A^*$ induces an isomorphism $\pi$ of $\A$.  As noted $\pi$ is induced by scalar multiplication by a rational $s$, there is an $n\in N$ such that for all $x$ and $y$, $\pi(x)=y$ if and only if $\A\models R(n,x,y)$.  Consequently, for $x^*$ and $y^*$ in $\A^*$, $\pi^*(x^*)=y^*$ if and only if $\pi(f(x^*))=f(y^*)$ if and only if $\A\models R(n,f(x^*),f(y^*))$ if and only if $\A^*\models R(n,x^*,y^*)$.  So, $\pi^*$ is defined in $A^*$ by a quantifier-free formula relative to the parameter $n$ and thereby computable from $X$.  It follows that  $\A$ is computably $\AUT$-countable on a cone, indeed on the cone above $\emptyset$.

\paragraph{Contradicting existential definability of automorphisms.}

The argument which follows is a generalization of the one that shows that $(\Q^*,\cdot)$ is not finitely generated.  To expedite the discussion, we say that a prime number $p$ is free over a finite set $F$ of rational numbers if and only if for every element of $x$ of $F$, if $a/b$ is the reduced-form fraction that represents $x$ then $p$ is not a factor of either $a$ or $b$.  We use the notation $pF$ to denote the set $\{px: x\in F\}$.

Suppose $\bar{c}=\{c_1,\dots,c_k\}$  is a finite subset of $Q$ and $\bar{n}=\{1,\dots,n\}$ is a finite initial segment of $N$.  We may assume that $0$ is an element of $\bar{c}$.  We propose to show that there is an automorphism $\pi$ of $\A$, given by multiplication by a non-zero rational number, such that $\pi$ is not definable by a countable family of existential formulas in the parameters $\bar{n}$, $\bar{c}$ and the images of $\bar{c}$ under $\pi$.  We will show this in a strong form.  Namely, we will show that if $p$ and $q$ are sufficiently free distinct prime numbers  and $\phi$ is an existential formula such that $\phi(\bar{n},\bar{c},p\bar{c},q,pq)$ holds in $\A$, then there is a number $y$ which is not equal to $pq$ such that $\phi(\bar{n},\bar{c},p\bar{c},q,y)$ also holds in $\A$.  The number $y$ will be obtained by multiplying $pq$ by a prime which is free over everything that came before. Note that this is sufficient to prove Theorem~\ref{3.1}.

With these comments in mind, let $\phi$ be  an existential formula in the language of $\A$.  Let $p$ be a prime number which is free over $\bar{c}$ and over all the $s_j\in S_m,$ for $m\in\bar{n}.$ We examine the automorphism of $\A$ given by multiplication within $Q$ by $p$. Let $q$ be a prime which is free over  $p$, $\bar{c},$  and all the $s_j\in S_m,$ for $m\in\bar{n}.$  Suppose that $\A\models\phi(\bar{n},\bar{c}, p\bar{c},q,pq)$, ostensibly to witness in $\A$ that multiplication by $p$ maps $q$ to $pq$.

For $\A$ to satisfy $\phi(\bar{n},\bar{c}, p\bar{c}, q,pq)$ there must be additional elements of $\A$ to witness the existential quantifiers in $\phi$, denoted by $\bar{n}^*=\{m^*_1,\dots,m^*_{\ell^*}\}\subset N\setminus \bar{n}$ and by $\bar{c}^*=\{c^*_1,\dots,c^*_{k^*}\}$ in $Q\setminus(\bar{c}\,\cup\,p\bar{c}\,\cup\,\{q,pq\}),$ 
such that $\A$ satisfies a conjunction of instances of $<$ and $R$ or their negations about elements of $\bar{c},p\bar{c},\bar{n},\bar{c}^*,\bar{n}^*, q$ and $pq$.  Fix such $\bar{c}^*$ and $\bar{n}^*$ and let $D$ be the set with elements $\bar{n},\bar{c},p\bar{c},\bar{n}^*,\bar{c}^*,q$ and $pq$.   Since there are only finitely many possibilities, without loss of generality, we may assume that the given conjunction includes all instances of the relations $<$ or $R$ or of their negations which hold of elements of $D.$   We call this the $\A$-diagram of $D$.

\newcommand{\Cnpq}{C(\bar{n},pq)}

By fixing the elements of $D$, we have also fixed all the ratios between the nonzero elements of $D\cap Q$, which determines not only  which of their triples satisfy $R$ but also by which multiplicative scalars.  Specifically, when $x$ is non-zero, $R(m,x,y)$ holds if and only if $y/x\in S_m$.  Consider  the $\bar{n}$-graph on $D\,\cap\,Q$ in which two points $x$ and $y$ are connected by an edge if and only if there is a $m\in\bar{n}$ such that either both $x$ and $y$ are zero or one of $x/y$ or $y/x$ is an element of $S_{m}$.   The $\bar{n}$-component of an element of $x\in D\,\cap\,Q$ is the set of $y\in D\,\cap\,Q$ such that there is a path in the $\bar{n}$-graph which goes from $x$ to $y$.  Let $\Cnpq$ denote the $\bar{n}$-component of $pq$.  

We observe that $\Cnpq$ does not include $q$ or any element of $\bar{c}\,\cup\,p\bar{c}$.  This is because multiplication of a rational number with reduced-form fractional representation $a/b$ by an element of any $S_{m}$, $m\in\bar{n}$, can neither add nor delete a factor of $p$ or $q$ from $a$ or $b$ in the fractional representation of the product. Neither $q$ nor any element of $\bar{c}\,\cup\,p\bar{c}$ is denoted by a fraction $a/b$ in reduced such that both $q$ and $p$ divide $a$.  But, by the same reasoning, all of the elements of $\Cnpq$ are denoted by such reduced-form fractions.

We exhibit a partially stretched copy $D^\dagger$ of $D$.  Let $r$ be a prime which is free over $q$, $p$, $\bar{c}$, all the $s_j\in S_m,$ for $m\in\bar{n},$ and all other elements of $D$.   First, replace $pq$ by $pqr$ and replace each element $c_i^*$ of $\Cnpq$ with $c^*_ir$.  In particular, do not replace $q$ or any element of $\bar{c}\,\cup\,p\bar{c}.$  Second, let $\bar{n}^{\dagger}$ be a sequence of indices which are pairwise ordered under $<$ in the same way as $\bar{n}^*$ is ordered and such that for each $m^*\in\bar{n}^*$, 
$S_{m^{\dagger}}$ is the smallest superset of $S_{m^*}$ such that the following conditions hold.
  \begin{enumerate}
  \item For each $x$ and $y$ in $D$, if $x$ is not in $\Cnpq$, $y$ is in $\Cnpq$ and $R(m^*,x,y)$ then $ry/x\in S_{m^{\dagger}}$.
  \item For each $x$ and $y$ in $D$, if $x$ is in $\Cnpq$, $y$ is not in $\Cnpq$ and $R(m^*,x,y)$ then $y/xr\in S_{m^{\dagger}}$.
    \end{enumerate}
Since each finite set $S$ appears infinitely often as an $S_j$ and all of the elements of $\bar{n}^*$ are larger than the maximum of $\bar{n}$, there is such a sequence. Replace each $m^*$ in $D$ by $m^{\dagger}$ in $D^\dagger$.
   
 
Thus, we have a map $\dagger:D\to D^\dagger$, which replaces the $\bar{n}$-component of $pq$ with a new component rescaled by multiplication by $r$ and which augments the sets $S_{m^*}$ to include the graphs of appropriately rescaled maps.  Intuitively, we have stretched $\Cnpq$ by a factor of $r$ and correspondingly stretched or contracted its connections to the rest of $D$, depending on their directions.

By inspection, $\dagger$ preserves membership in $N$ and membership in $Q$.

By the definition of  $\dagger$, for all $m_1^*$ and $m_2^*$ in $\bar{n}^*$,  $m^*_1<m^*_2$ if and only if $m^{\dagger}_1<m^{\dagger}_2.$  So, the relation $<$ is preserved by $\dagger$.

We examine all the cases to show the same for the relation $R$.  Consider a triple $(m,x,y)$.  Begin with the case $m\in\bar{n}$.  Then, $m^\dagger$ is equal to $m$.  There are three possible cases for pairs $(x,y)$ in $D^\dagger$: (1) neither comes from $\Cnpq$,   (2) one does and the other does not, or (3) both do.    In Case (1), $x^\dagger=x$, $y^\dagger=y$ and so $R(m,x,y)$ if and only if $R(m^\dagger,x^\dagger,y^\dagger)$ is true on trivial grounds.  In (2), $R(m,x,y)$ cannot hold, by the definition of $\Cnpq$.  We observe that neither $R(m,x,ry)$ nor $R(m,rx,y)$ can be true statements about elements of $D^\dagger$.  First, if $x\not\in\Cnpq$, $y\in\Cnpq$  and $s\in S_m$, then $r$ is free over $sx$ but not free over $ry$, so $R(m,x,ry),$ equivalently $R(m,x^\dagger,y^\dagger),$
also cannot hold.  The dual case, in which $x\in\Cnpq$ and $y\not\in\Cnpq,$ is similar.  In (3),  both $x$ and $y$ belong to $\Cnpq$.
Then $x^\dagger=rx$ and $y^\dagger=ry$ and $R(m,x,y)$ is equivalent to $R(m,x^\dagger,y^\dagger),$ by the observation that $R$ is invariant under scalar multiplication on the second and third coordinates.  

Now, we examine the case $m^*\in\bar{n}^*$.  By definition, $S_{m^\dagger}$ is equal to
\[
  \begin{gathered}
  S_{m^*}\,\cup\, \left\{
    sr: \exists x,y\in D\left(
      \begin{gathered}
        y\in\bar{n}\mbox{-component of }pq\mbox{ and }\\
        x\not\in\bar{n}\mbox{-component of }pq\mbox{ and }\\
        s=y/x\in S_{m^*} 
      \end{gathered}
    \right)\right\}\\
  \cup\left\{
    s/r: \exists x,y\in D\left(
      \begin{gathered}
        y\in\bar{n}\mbox{-component of }pq\mbox{ and }\\
        x\not\in\bar{n}\mbox{-component of }pq\mbox{ and }\\
        s=x/y\in S_{m^*} 
      \end{gathered}
    \right)\right\}
  \end{gathered}
  \]

  Again, consider the three cases for $x$ and $y$.  For Case (1), suppose that neither $x$ nor $y$ belongs to $\Cnpq$.  Then $x^\dagger=x$ and $y^\dagger=y$.  For any $m^*$, since $S_{m^*}$ is contained in $S_{m^\dagger}$, if $R(m^*,x,y)$ then $R(m^\dagger,x,y)$, which is the same as $R(m^\dagger,x^\dagger,y^\dagger)$.  Since $r$ is free over  $x$ and over $y$  and the additional elements of $R(m^\dagger,x,y)$ involve multiplication by a rational number with a factor of $r$ in either its denominator or in its numerator, if $R(m^\dagger,x^\dagger,y^\dagger),$ equivalently $R(m^\dagger,x,y)$, then $y/x\in S_{m^*}$ and $R(m^*,x,y).$ 

For (2), suppose  $y$ is in $\Cnpq$ and $x$ is not.  By definition of $\dagger$, $x^\dagger=x$ and $y^\dagger=ry$.  Suppose $R(m^*,x,y)$.   Then $y/x\in S_{m^*}$, which implies $ry/x\in S_{m^\dagger},$ which implies $R(m^\dagger,x,ry)$ , and so $R(m^\dagger,x^\dagger,y^\dagger)$.  Conversely, suppose that  $R(m^\dagger,x^\dagger,y^\dagger),$ equivalently $R(m^\dagger,x,ry)$.  By the definition of $R$,  $ry/x\in S_{m^\dagger}$.  Since $r$ is free over $x$, $y$ and $S_m$, $ry/x$ must have the form $rs$, for $s$ in $S_{m^*}$, and so $y/x\in S_{m^*}$, equivalently $R(m^*,x,y)$.  The case of $R(m^*,y,x)$ is similar, with division by $r$ in place of multiplication by $r$.  

Finally, suppose that both $x$ and $y$ belong to $\Cnpq.$  Then $x^\dagger=rx$ and $y^\dagger=ry$ and neither is equal to zero.  If $R(m^*,x,y)$, equivalently $y/x\in S_{m^*}$, then $ry/rx\in S_{m^*},$ equivalently $y^\dagger/x^\dagger\in S_{m^*}$.  Since $S_{m^*}\subseteq S_{m^\dagger}$, $y^\dagger/x^\dagger\in S_{m^\dagger}$ and so $R(m^\dagger,x^\dagger,y^\dagger)$ holds.    Conversely, suppose $R(m^\dagger,x^\dagger,y^\dagger)$.  Then  $y^\dagger/x^\dagger\in S_{m^\dagger}$, which is the same as saying $ry/rx\in S_{m^\dagger}$ and equivalent to $y/x\in S_{m^\dagger}.$   Since $r$ is free over $x$ and $y$, $y/x$ must be in $S_{m^\dagger}$ by reason of $y/x\in S_{m^*}$.  Thus, $R(m^*,x,y)$ must hold.  This exhausts all  the cases and proves that $R(m^*,x,y)$ holds if and only if $R(m^\dagger,x^\dagger,y^\dagger)$ also holds.

  Thus, $\dagger$ is an isomorphism between the relational structures on $D$ and $D^\dagger$.  Quantifier-free formulas are preserved by such maps, so  $\A\models\phi(\bar{n},\bar{c},p\bar{c},q,pqr)$, as was required to finish the proof of the theorem.
\end{proof}
  



\begin{thebibliography}{1}

\bibitem{AshKniManSla89}
Chris Ash, Julia Knight, Mark Manasse, and Theodore Slaman.
\newblock Generic copies of countable structures.
\newblock {\em Ann. Pure Appl. Logic}, 42(3):195--205, 1989.

\bibitem{Hod1993}
Wilfrid Hodges.
\newblock {\em Model theory}, volume~42 of {\em Encyclopedia of Mathematics and
  its Applications}.
\newblock Cambridge University Press, Cambridge, 1993.

\end{thebibliography}


\end{document}
